% Source: https://artofproblemsolving.com/community/c3h1887149_geometry
% Stable identity: aops-topic-1887149-q01
The diameter $AB$ is perpendicular to the chord $CD$, so it is the perpendicular bisector of $CD$ and $AC=AD=2\sqrt3$. Since $AB$ is a diameter, $\angle ACB=90^\circ$, so
\[
AB^2=AC^2+BC^2=12+36=48,\qquad AB=4\sqrt3,\qquad R=2\sqrt3 .
\]
Put the centre at the origin with $A=(-2\sqrt3,0)$ and $B=(2\sqrt3,0)$, and let $C=(h,k)$ with $k>0$. From $AC^2=12$ and $BC^2=36$,
\[
(h-2\sqrt3)^2-(h+2\sqrt3)^2=24\ \Longrightarrow\ -8\sqrt3\,h=24\ \Longrightarrow\ h=-\sqrt3,
\]
and $k^2=12-(h+2\sqrt3)^2=12-3=9$, so $C=(-\sqrt3,3)$ and $D=(-\sqrt3,-3)$.

\begin{asy}
usepackage("amsmath");
size(6cm);
defaultpen(fontsize(10pt));
real R=2*sqrt(3);
pair O=(0,0),A=(-R,0),B=(R,0),C=(-sqrt(3),3),D=(-sqrt(3),-3),H=(-sqrt(3),0);
pair N=(A+D)/2,M=(C+B)/2;
draw(circle(O,R));
draw(A--B);draw(C--D);draw(A--D);draw(C--B);
draw(M--N,red+linewidth(1));
draw(H+(0.3,0)--H+(0.3,0.3)--H+(0,0.3));
dot(O);dot(M);dot(N);
label("$A$",A,W);label("$B$",B,E);label("$C$",C,dir(C));label("$D$",D,dir(D));label("$O$",O,S);
label("$N$",N,dir(350));label("$M$",M,NE);
label("$6$",(-0.37,1.75));label("$2\sqrt3$",(-2.45,-0.62),fontsize(9pt));
label("\mbox{}",(min(currentpicture,user=true)+max(currentpicture,user=true))/2,black); // keeps the text after the figure black when LaTeX typesets the labels inline
\end{asy}

Then
\[
N=\dfrac{A+D}{2}=\left(-\dfrac{3\sqrt3}{2},-\dfrac{3}{2}\right),\qquad M=\dfrac{C+B}{2}=\left(\dfrac{\sqrt3}{2},\dfrac{3}{2}\right),
\]
and
\[
MN=\sqrt{\left(2\sqrt3\right)^2+3^2}=\boxed{\sqrt{21}}.
\]
